\documentclass{report}
\usepackage{amsfonts, amsmath}
\newcommand\R{{\mathbb R}}
\newcommand\Q{{\mathbb Q}}
\newcommand\N{{\mathbb N}}
\newcommand\ve{\varepsilon}
\newcommand\Co{{\mathcal C}}
\pagestyle{empty}
\begin{document}
\large
\centerline{\Huge Fisica Matematica II}
\renewcommand{\thefootnote}{ } 
\footnote{\sl Avete 2:30 ore di tempo. Ogni esercizio
vale dieci punti. La sufficienza si ottiene con un punteggio $\geq$
18. Solo le risposte chiaramente giustificate saranno prese in
considerazione.}
\renewcommand{\thefootnote}{\arabic{footnote}} 
\vskip.1cm
\centerline{\Large Appello 10-07-2003}
\vskip1cm
\begin{enumerate}
\item Si consideri la soluzione dell'equazione
\[
\begin{split}
&(y^2-u^2)u_x-xyu_y=xu\\
&u(s,s)=s\quad s>0.
\end{split}
\]
\item  Sia $Q:=(0,1)^2\subset\R^2$ e sia $u:\bar Q\to\R$ una
soluzione, non identicamente nulla, di
\[
\begin{split}
&\Delta u=-\lambda u\\
&u|_{\partial Q}=0.
\end{split}
\]
Si dimostri che deve essere $\lambda >0$. Si cerchi di determinare il
$\lambda$ minimo possibile.
\item Per ogni $a>0$ sia $u_a:\R\times\R_+\to\R$ la soluzione di
\[
\begin{split}
&\partial_t u_a=\Delta u_a\\
&u_a(x,0)=a^{-\frac 12}e^{-\frac {x^2}a}.
\end{split}
\]
Si calcoli
\[
\lim_{a\to 0} u_a(x,1).
\]
\item Si determini la soluzione dell'equazione
\[
\begin{split}
&u_{tt}=u_{xx}+e^{-t} \quad\quad 0<x<1;\;t>0\\
&u(x,0)=0;\; u_t(x,0)=0\quad  0<x<1\\
&u(0,t)=u(1,t)=0\quad t>0.
\end{split}
\]
Si calcoli inoltre $\lim_{t\to\infty} E(t)$ per
\[
E(t):=\frac 12\int_0^1 u_t^2+u_x^2\; dx.
\]
\end{enumerate}
\end{document}
